\documentclass[a4paper,12pt]{article}
\usepackage[T2A]{fontenc}
\usepackage[utf8]{inputenc}
\usepackage[russian]{babel}
\usepackage{amsmath}
\usepackage{graphicx}
\usepackage{caption}

\begin{document}

\section*{Детерминизм}

\begin{figure}[h!]
    \centering
    \includegraphics[width=0.8\textwidth]{1.png} 
    \caption{Стековый автомат}
    \label{fig:stack_automaton}
\end{figure}

\textbf{Накачка:}

\centerline{$b^{2n}aaab^{2n-1}$}
\vspace{0.3cm}
\centerline{$\swarrow$\hspace{2cm}$\searrow$}
\centerline{$b$\hspace{2.5cm}$bbaaab^{4n}$}

\medskip

\textbf{Пояснение:}
\begin{itemize}
\item Из рисунка видно, что только лишь в 13 состоянии имеется недетерминизм, докажем его.
    \item Если качать левое слово синхронно, то очевидно выпадем, так как стек не кончится.
    \item Если префикс качаем только, то нулевая накачка его выведет правое слово из языка, так как опять стек не кончится.
\end{itemize}

% Здесь будут дополненные главы

\section*{Беспрефиксность и LL-свойство}

Рассмотрим слово $w = bbaaabb$, которое лежит в языке. \\
Добавим к нему $v = abb$, получим $wv = bbaaabbabb$, которое также принадлежит языку.

Следовательно, язык \textbf{не обладает префиксным свойством}, так как существует слово $w$, являющееся префиксом другого слова $wv$ из языка, при этом само $w$ также принадлежит языку.

Также в предыдущем пункте было показано, что язык \textbf{недетерминирован}. \\
Поскольку LL-свойство требует детерминированности, делаем вывод: \\
\textbf{Язык не обладает LL-свойством}.
\section*{Построение аппроксимаций}

В начале заметим, что правило, по которому $bb$ и $aaa$ были получены из любого слова, определяется однозначно. То есть их можно получить строго по их правилу, нельзя получить их иначе (например, конкатенацией слов). Поэтому можно заменить их на $c$ и $d$ соответственно, чтобы упростить работу с автоматом (думается мне, так можно делать), затем в правилах вернём всё на место.

% Место для рисунка 1: позиционный автомат
\begin{figure}[h!]
    \centering
    \includegraphics[width=0.8\textwidth]{2.png}
    \caption{Позиционный автомат}
    \label{fig:positional_automaton}
\end{figure}

Получаем такой LR(0)-автомат. Преобразовав, получаем его эквивалент:

% Место для рисунка 2: преобразованный позиционный автомат
\begin{figure}[h!]
    \centering
    \includegraphics[width=0.8\textwidth]{3.png} 
    \caption{Преобразованный позиционный автомат}
    \label{fig:transformed_positional_automaton}
\end{figure}

С которого уже находим правила через пересечения.

\subsection*{Правила из LR(0)-аппроксимации}

\begin{align*}
&\langle 0, T, 3 \rangle \rightarrow \langle 0, T, 3 \rangle a \langle 1, T, 3 \rangle \\
&\langle 1, T, 3 \rangle \rightarrow \langle 1, T, 3 \rangle a \langle 1, T, 3 \rangle \\
&\langle 3, T, 3 \rangle \rightarrow \langle 3, T, 3 \rangle a \langle 1, T, 3 \rangle \\
&\langle 5, T, 6 \rangle \rightarrow \langle 5, T, 6 \rangle a \langle 4, T, 6 \rangle \\
&\langle 4, T, 6 \rangle \rightarrow \langle 4, T, 6 \rangle a \langle 4, T, 6 \rangle \\
&\langle 6, T, 6 \rangle \rightarrow \langle 6, T, 6 \rangle a \langle 4, T, 6 \rangle \\
&\langle 0, T, 3 \rangle \rightarrow bb \\
&\langle 1, T, 3 \rangle \rightarrow bb \\
&\langle 3, T, 3 \rangle \rightarrow bb \\
&\langle 5, T, 6 \rangle \rightarrow bb \\
&\langle 4, T, 6 \rangle \rightarrow bb \\
&\langle 6, T, 6 \rangle \rightarrow bb \\
&\langle 0, S, 6 \rangle \rightarrow \langle 0, T, 3 \rangle \langle 3, S, 5 \rangle \langle 5, T, 6 \rangle \\
&\langle 0, S, 6 \rangle \rightarrow \langle 0, T, 3 \rangle \langle 3, S, 6 \rangle \langle 6, T, 6 \rangle \\
&\langle 3, S, 6 \rangle \rightarrow \langle 3, T, 3 \rangle \langle 3, S, 5 \rangle \langle 5, T, 6 \rangle \\
&\langle 3, S, 6 \rangle \rightarrow \langle 3, T, 3 \rangle \langle 3, S, 6 \rangle \langle 6, T, 6 \rangle \\
&\langle 0, S, 5 \rangle \rightarrow \langle 0, S, 5 \rangle b \langle 0, S, 5 \rangle \\
&\langle 0, S, 5 \rangle \rightarrow \langle 0, S, 6 \rangle b \langle 0, S, 5 \rangle \\
&\langle 0, S, 6 \rangle \rightarrow \langle 0, S, 5 \rangle b \langle 0, S, 6 \rangle \\
&\langle 0, S, 6 \rangle \rightarrow \langle 0, S, 6 \rangle b \langle 0, S, 6 \rangle \\
&\langle 3, S, 5 \rangle \rightarrow \langle 3, S, 5 \rangle b \langle 0, S, 5 \rangle \\
&\langle 3, S, 5 \rangle \rightarrow \langle 3, S, 6 \rangle b \langle 0, S, 5 \rangle \\
&\langle 3, S, 6 \rangle \rightarrow \langle 3, S, 5 \rangle b \langle 0, S, 6 \rangle \\
&\langle 3, S, 6 \rangle \rightarrow \langle 3, S, 6 \rangle b \langle 0, S, 6 \rangle \\
&\langle 0, S, 5 \rangle \rightarrow aaa \\
&\langle 3, S, 5 \rangle \rightarrow aaa
\end{align*}

% Место для рисунка 3: LL(1) автомат
\begin{figure}[h!]
    \centering
    \includegraphics[width=0.8\textwidth]{4.png} 
    \caption{LL(1)-автомат}
    \label{fig:LL1_automaton}
\end{figure}

LL(1)-автомат тоже приведён на рисунке. Список получившихся с него правил:

\subsection*{Правила из LL(1)-аппроксимации}

\begin{align*}
&\langle 0, T, 1 \rangle \rightarrow \langle 0, T, 1 \rangle a \langle 4, T, 1 \rangle \\
&\langle 4, T, 1 \rangle \rightarrow \langle 4, T, 1 \rangle a \langle 4, T, 1 \rangle \\
&\langle 2, T, 1 \rangle \rightarrow \langle 2, T, 1 \rangle a \langle 4, T, 1 \rangle \\
&\langle 3, T, 1 \rangle \rightarrow \langle 3, T, 1 \rangle a \langle 4, T, 1 \rangle \\
&\langle 0, T, 1 \rangle \rightarrow bb \\
&\langle 4, T, 1 \rangle \rightarrow bb \\
&\langle 2, T, 1 \rangle \rightarrow bb \\
&\langle 3, T, 1 \rangle \rightarrow bb \\
&\langle 0, S, 1 \rangle \rightarrow \langle 0, T, 1 \rangle \langle 1, S, 3 \rangle \langle 3, T, 1 \rangle \\
&\langle 0, S, 1 \rangle \rightarrow \langle 0, T, 1 \rangle \langle 1, S, 2 \rangle \langle 2, T, 1 \rangle \\
&\langle 3, S, 1 \rangle \rightarrow \langle 3, T, 1 \rangle \langle 1, S, 3 \rangle \langle 3, T, 1 \rangle \\
&\langle 3, S, 1 \rangle \rightarrow \langle 3, T, 1 \rangle \langle 1, S, 2 \rangle \langle 2, T, 1 \rangle \\
&\langle 0, S, 1 \rangle \rightarrow \langle 0, S, 2 \rangle b \langle 3, S, 1 \rangle \\
&\langle 0, S, 1 \rangle \rightarrow \langle 0, S, 1 \rangle b \langle 3, S, 1 \rangle \\
&\langle 0, S, 2 \rangle \rightarrow \langle 0, S, 2 \rangle b \langle 3, S, 2 \rangle \\
&\langle 0, S, 2 \rangle \rightarrow \langle 0, S, 1 \rangle b \langle 3, S, 2 \rangle \\
&\langle 3, S, 1 \rangle \rightarrow \langle 3, S, 2 \rangle b \langle 3, S, 1 \rangle \\
&\langle 3, S, 1 \rangle \rightarrow \langle 3, S, 1 \rangle b \langle 3, S, 1 \rangle \\
&\langle 3, S, 2 \rangle \rightarrow \langle 3, S, 2 \rangle b \langle 3, S, 2 \rangle \\
&\langle 3, S, 2 \rangle \rightarrow \langle 3, S, 1 \rangle b \langle 3, S, 2 \rangle \\
&\langle 1, S, 2 \rangle \rightarrow \langle 1, S, 2 \rangle b \langle 3, S, 2 \rangle \\
&\langle 1, S, 2 \rangle \rightarrow \langle 1, S, 1 \rangle b \langle 3, S, 2 \rangle \\
&\langle 1, S, 3 \rangle \rightarrow \langle 1, S, 2 \rangle b \langle 3, S, 3 \rangle \\
&\langle 1, S, 3 \rangle \rightarrow \langle 1, S, 1 \rangle b \langle 3, S, 3 \rangle \\
&\langle 1, S, 1 \rangle \rightarrow \langle 1, S, 2 \rangle b \langle 3, S, 1 \rangle \\
&\langle 1, S, 1 \rangle \rightarrow \langle 1, S, 1 \rangle b \langle 3, S, 1 \rangle \\
&\langle 3, S, 3 \rangle \rightarrow \langle 3, S, 2 \rangle b \langle 3, S, 3 \rangle \\
&\langle 3, S, 3 \rangle \rightarrow \langle 3, S, 1 \rangle b \langle 3, S, 3 \rangle \\
&\langle 0, S, 2 \rangle \rightarrow aaa \\
&\langle 1, S, 2 \rangle \rightarrow aaa \\
&\langle 3, S, 2 \rangle \rightarrow aaa
\end{align*}

\end{document}